% !TeX spellcheck = it_IT

\documentclass[format=169]{beamer}

\usepackage[utf8]{inputenc}
\usepackage[italian]{babel}
\usepackage{multicol}

\usetheme{metropolis}

\usepackage{mathtools}
\DeclarePairedDelimiter{\ceil}{\lceil}{\rceil}

\title{Esercitazioni di Informatica B}
\subtitle{Matrici}

\author{Stefano Cereda\\
	stefano.cereda@polimi.it
}
\date{23/10/2018}
\institute[PoliMi]{\vspace{0.5cm}\centering Politecnico di Milano \\ \vspace{0.2cm}
	\includegraphics[width=0.2\textwidth]{../logopolimi}}

\setbeamercovered{invisible}

\makeindex

\begin{document}
	\begin{frame}
	\maketitle
	\end{frame}

\begin{frame}<handout:0>{Ricerca e ordinamento}
Completare gli esercizi sugli array: ricerca per bisezione e ordinamento a bolle.
\end{frame}

\begin{frame}{Creazione di una matrice identità}
Scrivere un programma che crei e visualizzi una matrice identità di dimensioni note $N \times N$.

Si modifichi il programma per consentire all'utente di scegliere le dimensioni della matrice, con il vincolo che le dimensioni sia inferiori ad $N \times N$.
\end{frame}

\begin{frame}{Matrice in valore assoluto}
Scrivere un programma che, ricevuta in ingresso una matrice di interi $A$, calcoli e visualizzi una seconda matrice $B$ tale che:
\[ B_{ij} = |A_{ij}| \]
\end{frame}

\begin{frame}{Matrice trasposta}
Data una matrice quadrata intera $M$, di dimensioni $D \times D$, si scriva il codice per calcolare e visualizzare la trasposta. (le righe diventano colonne e viceversa).
\end{frame}

\begin{frame}{Rotazione di 90 gradi}
Data una matrice intera $A$, di dimensioni $R \times C$, la si ruoti in senso orario di 90 gradi. (Si noti che il risultato è diverso dalla trasposta).
\end{frame}

\begin{frame}{Da array a matrice}
Scrivere un programma che riceva un numero noto $N=R\cdot C$ di interi e li salvi in una matrice $M$ di dimensioni note $R \times C$. Si stampi poi la matrice in formato tabellare.
\end{frame}


\begin{frame}{Prodotto fra matrici}
Date due matrici intere $A, B$ si scriva un programma che ne calcoli il prodotto $C=A\times B$

Si ricorda che il prodotto di due matrici $A$ e $B$ (di dimensioni $m \times n, n \times p$) è una matrice $C$ (di dimensioni $m \times p$) i cui elementi $C_{ij}$ sono dati da:
\[ c_{ij} = \sum_{r=1}^n  a_{ir}b_{rj} \]
\end{frame}

\begin{frame}{Array e matrici}
Si scriva un programma che riceva in input un numero $dim<MAX\_V$ ($MAX\_V=100$) e chieda poi $dim$ interi in modo da popolare un vettore $v$. 

Il programma richiede all’utente altri due numeri, $r$ e $c$ rappresentanti le dimensioni di una matrice $m$.
Il programma deve usare i primi $r*c$ valori di $v$ per riempire la matrice $m$.

\begin{itemize}
	\item Se $r*c > dim$ il programma termina scrivendo un messaggio di errore.
	\item Se $r*c < dim$ il programma riempie $m$ e stampa i valori di $v$ che sono stati esclusi
	dalla matrice.
	\item Se $r*c = dim$ il programma stampa un messaggio di successo.
\end{itemize}

Nei casi in cui la matrice viene popolata si ricerchi all’interno di essa il massimo e minimo e se ne
stampino i valori.
\end{frame}

\begin{frame}{Stringhe e matrici}
Si definisca la struttura dati composta da una matrice di stringhe (NxN).

Si scriva quindi il codice che permette ad un utente di riempire la matrice con parole.
Alla fine della fase di inserimento si deve stampare la lunghezza media delle parole ed il numero
totale di caratteri presenti nella matrice.

Si costruisca un vettore contenente le parole presenti nella diagonale della matrice e se ne stampi
il contenuto.

Si chiede di inserire una parola da ricercare fra il contenuto della matrice. Il risultato della ricerca restituisce gli indici ([riga,colonna]) della prima occorrenza della parola (se essa esiste) oppure il valore -1.
\end{frame}

\end{document}